% Abstract: 150-220 words (median 180). Sentence roles, in order: context; problem
% and stakes; why current approaches fail; "we present"; key idea; techniques
% paired with challenges; evaluation basis; headline numbers WITH baseline; code.
% Every number below must appear identically in the body (paper-state.md, "Numbers").
% Draft status: first draft, written early so the paper has a target. Rewrite in
% Step 5.6, after the evaluation is final.
Learned policies now beat hand-written heuristics in CPU caches, content-delivery
networks and storage caches, but an operating system's page-replacement policy works
under constraints those settings lack: the kernel never sees individual memory
accesses, cannot use floating point on the eviction path, and every eviction changes
all later faults.
We ask whether a learned page-replacement policy is implementable under these
constraints, and answer with \sys, built in the xv6 teaching kernel.
\sys scores resident pages with an integer linear model over signals the kernel
already keeps (accessed- and dirty-bit scan history and refault bookkeeping), is
trained offline to imitate \belady's optimal policy on recorded eviction decisions,
and shields newly loaded pages with a probation rule.
On 89 verified reference streams (123M references) from five workload families,
these kernel-observable features match oracle features to within 1.5\% on three
families, and 8-bit integer weights change fault counts by at most 0.002 of \clock's.
Running inside the kernel on workload variants never seen in training, \sys incurs
3--46\% fewer page faults than \clock and up to 96\% fewer page writes.
The price is decision time: choosing a victim costs 11--13$\times$ as much as \clock.
\todo{anonymized code link; update dataset size if the 7 new families enter the study}
